\section{Oil Poured Out: The Feast of the Most Holy Name of Jesus}\label{sec:holy-name}

In the closing clause of the Gospel of 1~January, ``his name was called JESUS,
which was called by the angel before he was conceived in the womb''
(Luke~2:21). Immediately after the Octave Day closes, the 1962 Roman calendar
sets apart an entire second-class Feast of the Lord to adore that single word:
the \textbf{Feast of the Most Holy Name of Jesus} (\textit{Sanctissimi Nominis
Iesu}).

\subsection{How the Feast Entered Christmastide and When It Falls}

Until the twentieth century, the Roman Missal and Breviary kept the Feast of
the Most Holy Name on the Second Sunday after the Epiphany (as seen in the 1862
Ratisbon Missal and the 1888 Pustet Breviary). On 23~October~1913, through his
motu proprio \textit{Abhinc duos annos} and the Sacred Congregation of Rites
decree of 28~October~1913, Pope St.~Pius~X permanently moved the feast into
Christmastide, assigning it to the \textbf{Sunday falling between 2 and
5~January}, or to \textbf{2~January} whenever no Sunday falls in that four-day
span.

The 1960 \textit{Codex Rubricarum} retained St.~Pius~X's exact rule in General
Rubric~17~a, ranking the feast at row~14 of the table of precedence (second-class
feast of the Lord). Across the century from 2020 through 2120, the Most Holy
Name falls on a Sunday in \textbf{58~years} (2, 3, 4, or 5~January) and on the
\textbf{2~January weekday fallback} in \textbf{43~years} (Monday, Tuesday, or
Wednesday, whenever 1~January falls on Sunday, Monday, or Tuesday). Under General Rubrics~16~a and~111~b, when the Most Holy Name
falls on \textbf{Sunday 5~January} (fourteen times per century, as in~2025 and
2031), it assumes the Sunday's full rights and privileges and \textbf{omits the
commemoration of Pope St.~Telesphorus entirely}, even at Low Mass.

At Mass, the vestments are white, the \textit{Gloria} and Creed are sung, and
the \textbf{Preface of the Nativity} continues to be used as the seasonal
preface (\textit{tamquam de Tempore}). Because the Christmas Octave ended on
1~January, however, the Canon returns to the ordinary \textit{Communicantes}.

\subsection{At the Name of Jesus Every Knee Shall Bow}

The Mass of the Most Holy Name welds St.~Paul, St.~Peter, and St.~Luke into one
confession of faith. The Introit takes up the great Christological hymn of
Philippians~2:10--11 joined to Psalm~8:2: \textit{In nomine Jesu omne genu
flectatur, caelestium, terrestrium, et infernorum: et omnis lingua confiteatur,
quia Dominus Jesus Christus in gloria est Dei Patris}---``In the name of Jesus
every knee should bow, of those that are in heaven, on earth, and under the
earth: and every tongue should confess that the Lord Jesus Christ is in the
glory of God the Father.'' The Collect asks that the worship of his Name on
earth may ripen into the beatific vision in heaven:

\begin{massbox}{Collect of the Most Holy Name of Jesus}
\textbf{Latin} (1962 Roman Missal). Deus, qui unigenitum Filium tuum
constituisti humani generis Salvatorem, et Jesum vocari jussisti: concede
propitius; ut, cujus sanctum nomen veneramur in terris, ejus quoque aspectu
perfruamur in caelis. Per eundem Dominum.

\textbf{English} (Bute, 1908). O GOD, Who hast appointed Thine Only-begotten
Son to be the Saviour of mankind, and hast commanded that His Name should be
called JESUS, mercifully grant that we who here on earth do worship that His
most Holy Name may be made glad in heaven by His Presence. Through the Same our
Lord Jesus Christ, Thy Son, Who liveth and reigneth with Thee, in the unity of
the Holy Ghost, One God, world without end. Amen.
\end{massbox}

Instead of an Epistle, the Mass reads St.~Peter's bold testimony before the
Sanhedrin after healing the lame man at the Beautiful Gate (Acts~4:8--12): by
the name of Jesus Christ of Nazareth, whom they crucified and God raised from
the dead, this man stands whole, for ``Neither is there salvation in any other.
For there is no other name under heaven given to men, whereby we must be
saved.'' The Gospel returns to the single verse of Luke~2:21---the circumcision
and naming of the Child on the eighth day---while the Offertory and Communion
draw from Psalm~85~[86]: ``All the nations thou hast made shall come and adore
before thee, O Lord: and they shall glorify thy name.''

In the \textit{Summa Theologiae} (Tertia Pars, q.~37, a.~2), St.~Thomas Aquinas
observes that names given by God in Scripture always signify a gratuitous gift
bestowed on the bearer: because the grace of universal salvation belongs to
Christ, the angel revealed the name \textit{Jesus} (``Saviour'') both to his
Mother and to St.~Joseph, gathering into one syllable every prophetic title---Emmanuel,
Wonderful, Counsellor, Orient. Just as Israel named children on the day of
circumcision, Thomas adds, Christians receive their names in Baptism, sealed
under the Name that saves.

\subsection{Honey in the Mouth, Music in the Ear, Gladness in the Heart}

In the Divine Office, the Feast of the Most Holy Name breathes the warmth of
Cistercian devotion. Lauds opens with Canticles~1:2, \textit{Oleum effusum
nomen tuum}, ``Thy name is as oil poured out,'' and the three nocturns of
Matins place St.~Bernard of Clairvaux's famous Fifteenth Sermon on the Song of
Songs (Lessons~4--6) alongside his First Sermon on the Circumcision
(Lessons~7--9).

In Lesson~4--6, St.~Bernard unpacks why the Holy Spirit compares the Name of
Jesus to poured-out oil:

\begin{quote}
Oil indeed giveth light, meat, and unction. It feedeth fire, it nourisheth the
flesh, it sootheth pain; it is light, food, and healing. Behold, thus also is
the Name of the Bridegroom. To preach it, is to give light; to think of it, is
to feed the soul; to call on it, is to win grace and unction... If thou writest,
thou hast no meaning for me if I read not of JESUS there. If thou preach, or
dispute, thou hast no meaning for me if I hear not of JESUS there. The mention
of JESUS is honey in the mouth, music in the ear, and gladness in the heart. It
is our healing too. Is any sorrowful among us? Let the thought of Jesus come
into his heart, and spring to his mouth. Behold, when the day of that Name
beginneth to break, every cloud will flee away, and there will be a great calm.
\end{quote}

In his Second Sermon for the Circumcision, Bernard deepens the thought: ``It
was called, but not given. For Jesus has been His name from everlasting. To be
a Saviour is a property of His Nature,'' and every other prophetic title holds
the Name of Jesus ``in a state of solution.'' Six centuries later,
St.~Alphonsus Maria de Liguori echoed Bernard in his \textit{Discourse on the
Name of Jesus}, marveling that at the very moment the Child humbled himself to
bear the mark of a sinner in his circumcision, the Father honored him with a
name above every name---``for it is this name that consoles us, defends us, and
makes us burn with love.''

Across Vespers, Matins, and Lauds, the Breviary sings three companion hymns
drawn from the medieval poem \textit{Jesu dulcis memoria}---at Vespers,
\textit{Jesu dulcis memoria}, ``Jesus! the very thought of Thee / With
sweetness fills the breast!''; at Matins, \textit{Jesu Rex admirabilis}, ``O
Jesu, King most wonderful!''; and at Lauds, \textit{Jesu decus angelicum}, ``O
Jesus, Thou the Beauty art / Of angel worlds above!'' (Appendix~\ref{app:prayers}).
Few devotions fit the busy days of early January more naturally than the one
this feast teaches: bowing the head slightly whenever the Name of Jesus is
spoken at Mass or in prayer, and calling quietly upon his Name in temptation,
fatigue, or sorrow, until ``every cloud will flee away, and there will be a
great calm.''
